% =============================================================
%  HeInSulfides_SI_body.tex
%  Shared, class-agnostic body of Chapter 3's Supplementary
%  Information. \input{} by Chapter3/chapter3_SI.tex, following the
%  exact wrapper/body split used by Chapter1/chapter1_SI.tex and
%  Chapter2/chapter2_SI.tex (see the comment in either for the
%  rationale). Edit this file, not chapter3_SI.tex, to change the
%  text, figures, tables, or citations.
%
%  Citations use \citep/\citet against the master references.bib,
%  same as HeInSulfides_Main_Body.tex -- resolved by the single
%  \printbibliography at the end of dissertation.tex.
%
%  Contents:
%    \S\ref{ch3-si-cu-pb-sulfides} -- the chalcopyrite/chalcocite/galena
%      step-heat Results and Discussion content moved here (at the
%      user's request) from the main text, which now only narrates
%      the pyrite/arsenopyrite results directly relevant to this
%      chapter's U-Th/He geochronology aims.
%    \S\ref{ch3-si-geochron} -- the complete U-Th/He geochronology
%      data compilation not discussed in the main text (Table~\ref{tab:ch3-si-geochron}),
%      plus the Gympie 4He-vs-age SI figure justifying the exclusion
%      of the 0413PyGy-4 outlier from the main-text figure. This
%      section is a data appendix, not a narrative -- future SI
%      material can be added here without needing to tell a story.
% =============================================================

\section{Chalcopyrite, Chalcocite, and Galena Step Heat Results and Discussion}
\label{ch3-si-cu-pb-sulfides}

This section reports the step heat results and discussion for the two
non-pyrite/arsenopyrite sulfides of this study -- the copper sulfides
chalcopyrite (CCpTS230630) and chalcocite (CCT230707), and galena
(GalTS230630) -- moved here from the main text because neither mineral is
a promising U-Th/He geochronometer or helium isotope target, and neither
is otherwise part of this chapter's main narrative (Table 1 in the main
text still lists all nine step heat experiments, including these three,
for completeness).

\subsection*{Results}

\subsubsection*{Chalcocite and Chalcopyrite}
Step heat data step heat experiments for chalcopyrite (CCpTS230630) and
chalcocite (CCT230707) produced macroscopic evidence of phase
transformation, including the development of an iridescent blue sulfide
or oxide layer on both specimens and dramatic crystal growth on the
ampoule walls for the chalcocite experiment \hl{Tube Photo} Although
permeation lag was not sufficiently constrained to calculate apparent
diffusivities for the chalcocite experiment, the data for CCT230707 are
consistent with a primary phase-breakdown control of helium migration at
temperatures $\leq$300°C.

A linear regression is fit to 9 measurements defining a linear region of
the Arrhenius plot of chalcopyrite sample CCpTS230630 for use as a
reference for the delta plot (Fig.~\ref{fig:ch3-si-stepheat-grid-other}, row 3). This regression has an
R\textsuperscript{2} value of \ChThreeCCpRtwo, an implied activation energy of
\ChThreeCCpEaKcal\ $\pm$ \ChThreeCCpEaKcalErr\ kcal (\ChThreeCCpEa\ $\pm$
\ChThreeCCpEaErr\ kJ/mol), and an implied closure temperature of
\ChThreeCCpTc\ $\pm$ \ChThreeCCpTcErr°C assuming a
cooling rate of 10 °C/Ma. The delta plot generated using this reference
line for chalcopyrite clearly shows a dramatic mid-experiment increase
in diffusivity following steps above 400°C, followed by a continuous
decrease in delta values (Fig.~\ref{fig:ch3-si-stepheat-grid-other}, row 3), similarly to that observed for
the dehydrating goethite sample Win-06-01b \citep{farley2024}.


\subsubsection*{Galena}
The majority of the helium total yield was evolved during the step heat
experiment of the encapsulated galena, in stark contrast to the
comparatively retentive pyrites and arsenopyrite. The first 70\% of helium
released in the galena experiment could not be reduced into diffusivities, because $\lambda t \leq 4$ and so 
steady state permeation out of the ampoule could not be assumed. There were no permeation steps for this experiment. 
These data are still included in Fig.~\ref{fig:ch3-fo-comparison}, because the cumulative yields are less affected by
the permeation lag, which mostly affects the short-term knowledge of helium degassing: since earlier degassed helium will have long since 
permeated out, allowing for generally accurate accounting of evolved helium throughout the experiment. 
The reference line from the chalcopyrite experiment showed a similar activation energy to
the activation energy implied by curved retrogrades of the galena
experiment (Fig.~\ref{fig:ch3-si-stepheat-grid-other}, row 2), so it was used as the reference for the delta
plot of galena. This delta plot shows a monotonic decrease in delta
values as the experiment proceeds (Fig.~\ref{fig:ch3-si-stepheat-grid-other}, row 2).

The full $^3$He step-heat datasets for both experiments, in the same
format as the main-text pyrite/arsenopyrite tables
(\S\ref{ch3-sulfide-step-heat-experiments}), are given in
Tables~\ref{tab:ch3-stepheat-gal} and~\ref{tab:ch3-stepheat-ccp}.

\emph{Sphalerite}
Although we conducted a step heat experiment on translucent sphalerite (SphTS230630), none of the sample's steps were above $\lambda t \leq 4$ and there were insufficient permeation
steps to constrain the permeation lag, so no diffusivities could be calculated. Qualitatively, the sphalerite experiment showed that it is somewhat less helium retentive than apatite, 
and it was consistent with a single diffusion domain, although our data do not permit any confident interpretation. 
Future work on sphalerite should conduct disciplined permeation steps to track the helium inventory of the ampoule 
following the procedure outlined in \S\ref{permeation-step-method}.

\ChThreeStepHeatGal
\ChThreeStepHeatCCp

\begin{figure}[p]
  \centering
  \includegraphics[width=\textwidth,height=\textheight,keepaspectratio]{Sulfide_StepHeat_SI_Grid_Pyrite.pdf}
  \caption{Compiled Arrhenius (left) and delta (right) plots for the four
  plotable Chapter 3 pyrite step-heat experiments: (1) Rustenburg pyrite,
  (2) LVCOPy230630 (Leadville, small), (3) 112123LVCOPyIrr (Leadville,
  sulfur-isolated), (4) BigLVCOPy230707 (Leadville, big). Each row's
  Arrhenius panel plots the reference regression used for that row's
  delta panel (solid where owned, dashed where borrowed; see
  \S\ref{ch3-sulfide-step-heat-experiments}). Each row has its own x-axis
  scale, labeled on every panel. The three non-pyrite sulfides are
  compiled separately in Fig.~\ref{fig:ch3-si-stepheat-grid-other}.}
  \label{fig:ch3-si-stepheat-grid-pyrite}
\end{figure}

\begin{figure}[p]
  \centering
  \includegraphics[width=\textwidth,height=\textheight,keepaspectratio]{Sulfide_StepHeat_SI_Grid_Other.pdf}
  \caption{Compiled Arrhenius (left) and delta (right) plots for the
  three plotable Chapter 3 non-pyrite sulfide step-heat experiments:
  (1) SalAspIrr (arsenopyrite), (2) GalTS230630 (galena),
  (3) CCpTS230630 (chalcopyrite). Each row's Arrhenius panel plots the
  reference regression used for that row's delta panel (solid where
  owned, dashed where borrowed; see
  \S\ref{ch3-sulfide-step-heat-experiments}). Each row has its own
  x-axis scale, labeled on every panel. The four pyrite samples are
  compiled separately in Fig.~\ref{fig:ch3-si-stepheat-grid-pyrite}.}
  \label{fig:ch3-si-stepheat-grid-other}
\end{figure}

\subsection*{Discussion: Thermal Helium Migration Systematics}

\emph{Chalcopyrite and Chalcocite}

Neither copper-bearing sulfide mineral studied here was successfully
stabilized by encapsulation with sulfur, so these results provide little
insight into helium migration within copper sulfides in nature, though they
suggest high potential for helium retention at surface temperatures, similar to apatite. Still,
they are unlikely to be promising targets for U-Th/He geochronology or
helium isotope studies as there exists many copper-bearing sulfide
phases stable at temperatures between 200-500°C, for which phase
stability and Cu/Fe ratios are sensitive functions of temperature,
composition, and sulfur fugacity \citep{zhao2014,chenharvey1975,yundkullerud1966}. These
reactions are likely to result in helium loss during reheating, burial, diagenesis, or metasomatism after crystallization.
Copper sulfides are unlikely to be promising targets for U-Th/He geochronology, but may be suitable for helium isotope studies in young or unaltered samples.

\emph{Galena}

Because the majority of the cumulative yield was degassed in steps for which $\lambda t > 4$, the bulk helium behavior for galena is somewhat unclear.
The data of the remaining 30\% of the cumulative yield was consistent with a multiple diffusion domain helium loss
mechanism. As expected for an MDD system, the delta plot for galena
monotonically decreases as the fractional yield increases (Fig.~\ref{fig:ch3-si-stepheat-grid-other}, row 2;~\citet{farley2018}).
The majority of the 3He cumulative fractional yield was evolved at low temperatures where the
uncertainties due to permeation lag could not be constrained, yet these early degassing domains were certainly less retentive than 150 µm
half-prism width Durango apatite \citep{farley2000}. Additionally, the total 3He concentration for the irradiated galena was less than half that of
co-irradiated pyrite, implying poor helium retention even at ambient temperatures.
Galena is unlikely to be a promising target for U-Th/He geochronology or helium isotope studies.

\section{Complete U-Th/He Geochronology Compilation}
\label{ch3-si-geochron}

\subsection*{Trace Elements for the Main-Text Sample Suites}

The main-text geochronology table (Table~\ref{tab:ch3-main-geochron}) reports
U, Th, Sm, and Zr for the Kidd Creek, Gympie, Durango, and Gogebic aliquots.
The remaining ICPMS trace elements measured the same way -- Cu, As, Ag, and
Sb -- are reported here instead (Table~\ref{tab:ch3-main-geochron-trace}),
in the same row order, to keep the main-text table within the page width.

\ChThreeGeochronMainTraceSITable

\subsection*{Additional Sample Suites}

This subsection is a data appendix, not a narrative: it reports the complete
U-Th/He geochronology compilation (Chapter3/Tables/geochronology\_
compilation\_BJT.xlsx) for the sample suites and replicate aliquots not
discussed in the main text -- Kidd Creek replicate/mass-only aliquots,
Franklin (NJ) pyrite and fluorite, Rustenberg/Pilanesberg pyrite, Magnet
Cove carbonatite pyrite, French Creek magnetite/pyrite, and Salzburg
arsenopyrite -- so that the full underlying dataset is available even
where it is not part of this chapter's argument. Later work extending
this chapter's geochronology (e.g. to these additional localities) can
build directly on this section.

\ChThreeGeochronSITable

\begin{figure}[htbp]
  \centering
  \includegraphics[width=0.7\textwidth]{Gympie_Age_vs_4He_SI.pdf}
  \caption{Companion to the main-text Gympie age-vs-eU figure
  (Fig.~\ref{fig:ch3-gympie-age-vs-eu}), plotting atmospherically
  corrected age against measured 4He content (nccSTP/g) for the
  newer-protocol Gympie pyrite suite with the 0413PyGy-4 outlier put back
  in (\S\ref{ch3-results-gympie-ages}, main text). Its 4He content stands
  well off the trend defined by the other eleven aliquots, consistent
  with the helium-implantation interpretation offered there.}
  \label{fig:ch3-gympie-4he-outlier-si}
\end{figure}
